\subsection{Dual of a quotient space and of a subspace}

Throughout this section, E denotes a locally convex space, M a vector subspace of E, and \(M^{\circ}\) the orthogonal of M in the dual \(E'\) of E. Let p be the canonical mapping from E onto E/M; then \(^{t}p\) is injective, with image \(M^{\circ}\) , hence defines a vector space isomorphism (not topological)

\[
\pi : (\mathrm{E} / \mathbf {M}) ^ {\prime} \rightarrow \mathbf {M} ^ {\circ}.
\]

Similarly, let i be the canonical injection from M into E. Then \({}^{t}i\) is surjective (II, p.~24, prop. 2); its kernel is equal to \(M^{\circ}\) , and we get a vector space isomorphism (not topological)

\[
\iota : \mathrm{E} ^ {\prime} / \mathrm{M} ^ {\circ} \rightarrow \mathrm{M} ^ {\prime}.
\]

PROPOSITION 9. --- (i) For a subset A of (E/M)' to be equicontinuous, it is necessary and sufficient that \(\pi(A)\) is an equicontinuous subset of \(E'\) .

\begin{enumerate}
\def\labelenumi{(\roman{enumi})}
\setcounter{enumi}{1}
\item
  Let \(\mathfrak{S}\) be a set of bounded subsets of E, and \(\mathfrak{S}_1\) the set of the images of subsets \(\mathrm{A} \in \mathfrak{S}\) in E/M. Then \(\pi\) is an isomorphism from \((\mathrm{E} / \mathrm{M})_{\mathfrak{S}_1}'\) onto \(\mathbf{M}^\circ\) , where \(\mathbf{M}^\circ\) is assigned the topology induced by that of \(\mathrm{E}_{\mathfrak{S}}'\) .
\item
  Suppose E is a normed space then \(\pi\) is an isometry from the normed space (E/M)' onto the normed subspace M° of E'.
\end{enumerate}

Let A be a subset of \((\mathrm{E}/\mathrm{M})'\) and \(\mathbf{B} = {}^{t}p(\mathbf{A}) \subset \mathbf{E}'\) . Put

\[
q (\xi) = \sup _ {\xi^ {\prime} \in \mathbf {A}} | \langle \xi , \xi^ {\prime} \rangle |
\]

for all \(\xi \in \mathrm{E} / \mathbf{M}\) . In order that A be equicontinuous, it is necessary and sufficient that the mapping \(q\) from \(\mathrm{E} / \mathbf{M}\) into \(\overline{\mathbf{R}}_+\) is a continuous semi-norm. This implies that \(q \circ p\) is a continuous semi-norm on \(\mathrm{E}\) (II, p.~27, prop. 5, (ii)). Since we have

\[
(q \circ p) (x) = \sup _ {x ^ {\prime} \in \mathbf {B}} | \langle x, x ^ {\prime} \rangle |
\]

for all \(x \in E\) , this in turn implies that B is equicontinuous in \(E'\) , and (i) follows. Let \(A \in S\) and let f be a continuous linear form on E/M. For every \(\lambda \in R_{+}\) , we have \(|f| \leqslant \lambda\) on \(p(A)\) if and only if \(|^{t}p(f)| \leqslant \lambda\) on A; hence (ii).

Finally we prove (iii). Let \(y'\) be in \((\mathrm{E} / \mathbf{M})'\) . An element in \(\mathrm{E} / \mathbf{M}\) has norm \(< 1\) if and only if it is the image under \(p\) of an element of norm \(< 1\) in \(\mathbf{E}\) . Hence

\[
\begin{array}{l} \| y ^ {\prime} \| = \sup _ {y \in \mathrm{E} / \mathbf {M}, \| y \| <   1} \left| \langle y, y ^ {\prime} \rangle \right| = \sup _ {x \in \mathrm{E}, \| x \| <   1} \left| \langle p (x), y ^ {\prime} \rangle \right| \\ = \sup _ {x \in \mathrm{E}, \| x \| <   1} \left| \langle x, ^ {t} p (y ^ {\prime}) \rangle \right| = \left\| ^ {t} p (y ^ {\prime}) \right\|, \end{array}
\]

and \({}^t p\) induces an isometry from (E/M)' onto \({\mathbf{M}}^{ \circ  }\) .

PROPOSITION 10. --- (i) For a subset A of M' to be equicontinuous, it is necessary and sufficient that it is the image under \({}^t i\) of an equicontinuous subset of E'.

\begin{enumerate}
\def\labelenumi{(\roman{enumi})}
\setcounter{enumi}{1}
\item
  Suppose M is closed in E. Let S be a covering of E consisting of bounded subsets and let \(S_{1}\) be the set of subsets of M of the form \(M \cap A\) for A in S. The bijective linear mapping \(\iota\) from \(E_{S}^{\prime}/M^{\circ}\) onto \(M_{S_{1}}^{\prime}\) is continuous. It is a homeomorphism if S is a directed set for the relation \(\subset\) and consists of closed convex and compact sets for \(\sigma(\mathrm{E}, \mathrm{E}')\) .
\item
  If E is assumed to be normed, then \(\iota\) is an isometry from \(E'/M^{\circ}\) onto \(M'\) .
\end{enumerate}

The image under \({}^t i\) of an equicontinuous subset of \(\mathbf{E}'\) is an equicontinuous subset of \(\mathbf{M}'\) (IV, p.~7, prop. 7). Conversely, let \(\mathbf{A}\) be an equicontinuous subset of \(\mathbf{M}'\) . The topology of \(\mathbf{M}\) is defined by the set of restrictions to \(\mathbf{M}\) of the continuous seminorms on \(\mathbf{E}\) . Hence there exists a continuous semi-norm \(p\) on \(\mathbf{E}\) such that \(|f(x)| \leqslant p(x)\) for all \(f \in \mathbf{A}\) and for all \(x \in \mathbf{M}\) . Let \(\mathbf{B}\) be the set of all linear forms \(g\) on \(\mathbf{E}\) such that \(|g| \leqslant p\) and whose restriction to \(\mathbf{M}\) belongs to \(\mathbf{A}\) . The set \(\mathbf{B}\) is equicontinuous in \(\mathbf{E}'\) ; by Hahn-Banach theorem (II, p.~23, cor. 1), we have \({}^t i(\mathbf{B}) = \mathbf{A}\) , hence (i) follows.

We now prove (ii). By prop. 6 of IV, p.~6, the linear mapping \({}^t i\) from \(\mathbf{E}_{\mathfrak{S}}^{\prime}\) into \(\mathbf{M}_{\mathfrak{S}_1}^{\prime}\) is continuous, and defines, by passing to the quotient, a continuous linear mapping \(\iota\) from \(E_{\mathfrak{S}}^{\prime}/M^{\circ}\) onto \(M_{\mathfrak{S}_{1}}^{\prime}\) . Let T be the topology on \(M^{\prime}\) obtained by transferring that of \(E_{\mathfrak{S}}^{\prime}/M^{\circ}\) by \(\iota\) ; this is finer than the \(S_{1}\) -topology.

Suppose now that \(\mathfrak{S}\) is a directed set for \(\subset\) and consists of closed, convex, balanced and compact sets for \(\sigma(\mathrm{E},\mathrm{E}')\) . To show that \(\iota\) is a homeomorphism, i.e.~that \(\mathcal{T}\) is coarser than the \(\mathfrak{S}_1\) -topology on \(\mathbf{M}'\) , it is enough to prove that \(\mathcal{T}\) is compatible with the duality between \(\mathbf{M}'\) and \(\mathbf{M}\) and that every equicontinuous set in \(\mathbf{M}\) (considered as the dual of \(\mathbf{M}\) with \(\mathcal{T}\) ) is contained in the homothetic of a set belonging to \(\mathfrak{S}_1\) . Since \(\mathcal{T}\) is finer than the \(\mathfrak{S}_1\) -topology and \(\mathfrak{S}_1\) is a covering of \(\mathbf{M}\) , the linear form \(y' \mapsto \langle y, y' \rangle\) on \(\mathbf{M}'\) is continuous for \(\mathcal{T}\) for every \(y \in \mathbf{M}\) . Let \(f\) be a linear form on \(\mathbf{M}'\) which is continuous for \(\mathcal{T}\) ; then \(f \circ {}^t i\) is a continuous linear form on \(\mathrm{E}_{\mathfrak{S}}'\) . The \(\mathfrak{S}\) -topology on \(\mathrm{E}'\) is coarser than the Mackey topology \(\tau(\mathrm{E}',\mathrm{E})\) ; for, the mapping \(d_{\mathrm{B}}: \mathrm{E} \to \mathrm{E}'^*\) is continuous for the topologies \(\sigma(\mathrm{E},\mathrm{E}')\) and \(\sigma(\mathrm{E}'^*,\mathrm{E}')\) , and since the latter is Hausdorff, the image under \(d_{\mathrm{B}}\) of a set which is compact for \(\sigma(\mathrm{E},\mathrm{E}')\) is compact for \(\sigma(\mathrm{E}'^*,\mathrm{E}')\) . By lemma 1 of IV, p.~3, there exists \(x_0 \in E\) such that \(f(^t i(x')) = \langle x_0, x' \rangle\) for all \(x' \in E'\) . In particular, \(\langle x_0, x' \rangle = 0\) for all \(x' \in M^\circ\) , and since \(\mathbf{M}\) is closed in \(\mathrm{E}\) , we have \(x_0 \in \mathbf{M}\) (II, p.~45, cor. 2); and finally, \(f(y') = \langle x_0, y' \rangle\) for all \(y' \in M'\) . This proves that \(\mathcal{T}\) is compatible with the duality between \(\mathbf{M}\) and \(\mathbf{M}'\) .

Now let A be a subset of M equicontinuous for the topology T on \(M'\) . By the definition of T, and in view of the hypothesis that S is directed, this means that there exists a set \(B \in S\) containing 0 and such that the upper bound \(\lambda\) of the numbers \(\left|\langle y, x'\rangle\right|\) for \(y \in A\) and \(x' \in B^{\circ}\) , is finite (III, p.~19, prop. 7). Since B is closed in E, the theorem of bipolars (II, p.~44, th. 1) shows that we have \(A \subset \lambda(B \cap M)\) ; this completes the proof of (ii).

We shall now prove (iii). Let \(y' \in M'\) . We shall prove the formula

\[
\| y ^ {\prime} \| = \inf _ {{}^t i (x ^ {\prime}) = y ^ {\prime}} \| x ^ {\prime} \|.\tag{6}
\]

By prop. 8, (ii) of IV, p.~7, we have \(\| {}^t i\| = \| i\|\) , and so \(\| {}^t i\| \leqslant 1\) , and

\[
\| y ^ {\prime} \| \leqslant \inf _ {{}^t i (x ^ {\prime}) = y ^ {\prime}} \| x ^ {\prime} \|.\tag{7}
\]

By Hahn-Banach theorem (II, p.~23, cor. 3), there exists a linear form \(x_0'\) on E which extends \(y'\) and is of the same norm; hence we get the inequality opposite to (7), since \({}^t i(x_0') = y'\) .

Remark. --- We know (II, p.~48, prop. 7, (ii)) that \(\iota\) is a topological vector space isomorphism from \(\mathrm{E}_s^{\prime}|\mathbf{M}^{\circ}\) onto \(\mathbf{M}_s^\prime\) (weak duals). For the topology of compact convex convergence, prop. 10 shows that \(\iota\) is an isomorphism from \(\mathrm{E}_{cc}^{\prime}/\mathrm{M}^{\circ}\) onto \(\mathbf{M}_{cc}^{\prime}\) when E is Hausdorff and M closed in E. For the strong topologies, \(\iota\) is a continuous mapping from \(\mathrm{E}_b^{\prime}/\mathrm{M}^{\circ}\) onto \(\mathbf{M}_b^{\prime}\) ; it is an isomorphism if E is a Banach space * or if E is semi-reflexive and M is closed in E (IV, p.~15) *, but this is not always so if E is a Fréchet space (IV, p.~58, exerc. 5, c)).

PROPOSITION 11. --- (i) The weakened topology on E/M is the quotient of that on E; the weakened topology on M is induced by that of E.

\begin{enumerate}
\def\labelenumi{(\roman{enumi})}
\setcounter{enumi}{1}
\tightlist
\item
  The Mackey topology on E/M is the quotient of that on E; the Mackey topology on M is finer than the topology induced by \(\tau(\mathrm{E}, \mathrm{E}')\) .
\end{enumerate}

Assertion (i) follows from prop. 7 of II, p.~48.

The canonical injection \(i:M\to E\) is continuous for the weakened topologies, hence for the Mackey topologies \(\tau(\mathbf{M},\mathbf{M}^{\prime})\) and \(\tau(\mathbf{E},\mathbf{E}^{\prime})\) (IV, p.~7, prop. 7). Similarly, the canonical projection \(p:E\to E/M\) is continuous for the Mackey topologies. We see immediately that the quotient topology on E/M obtained from \(\tau(\mathbf{E},\mathbf{E}^{\prime})\) is compatible with the duality between E/M and \((\mathrm{E}/\mathrm{M})'\) , hence is coarser than the Mackey topology on E/M, by Mackey's theorem (IV, p.~2, th. 1). This proves (ii).

